% !TEX program = pdflatex
% Paper outline — Replica vs.\ EC hybrid redundancy + MILP (SkyStore extension).
% Rhetorical structure and section flow modelled on IEEE PDF extract *_pdf_extract.txt*
% (``Fast Variable-Grained Resemblance Data Deduplication For Cloud Storage'').
%
% Build (from repo root): pdflatex -output-directory=paper/paper0511 paper/paper0511/main.tex
% Second pass if citations complain: bibtex paper/paper0511/main && pdflatex ...
%
\documentclass[conference]{IEEEtran}
\usepackage{cite}
\usepackage{amsmath,amssymb}
\usepackage{graphicx}
\usepackage{booktabs}
\usepackage{xspace}
\usepackage{url}

\hyphenation{SkyStore geo-distrib-uted}

\begin{document}

\title{Replica--Erasure Hybrid Redundancy for Geo-Distributed Cloud Storage:\\
MILP-Guided Minimum-Cost Placement Beyond SkyStore}

% Customize affiliation blocks before submission.
\author{\IEEEauthorblockN{Author Name(s)}
\IEEEauthorblockA{Department / Institution\\
City, Country\\
email@domain}}

\maketitle

\begin{abstract}
With the prevalence of multi-region cloud storage and trace-driven evaluation frameworks, \emph{SkyStore} provides a practical substrate for reasoning about replication and TTL policies under heterogeneous regional prices.
Nevertheless, the baseline redundancy storyline emphasizes full-object replicas and rule-driven placement, while \emph{erasure coding (EC)} is often treated orthogonally and seldom optimized jointly under a single \emph{minimum-cost} surrogate together with replica placement.
Currently, operators still struggle when egress and storage prices differ across regions: placing an optional second replica versus striping two coded shards yields a rapidly growing discrete decision space.
Moreover, purely heuristic tuning deteriorates cost effectiveness as the number of regions increases because feasible layouts cannot be enumerated at trace cadence.
Therefore, we propose an extension that introduces \textbf{hybrid replica--EC redundancy} and encodes periodic placement decisions as a compact \textbf{mixed-integer linear program (MILP)}, minimizing linearized storage and placement-transfer charges derived from a cloud-style topology graph (storage prices plus edge egress attributes).
Under fixed-base semantics the mandatory base replica is modeled outside replica indicators via an objective constant over PUT-to-horizon duration; optional replicas absorb TTL snapshots while EC activates two outbound data stripes (\(size/K_{\mathrm{data}}\)) without a separate parity stripe variable in the MILP.
Finally, we implement the MILP-guided policy atop a SkyStore-derived discrete-event simulator and conduct trace-driven experiments together with analytic illustrations; the results show that the minimum-cost objective systematically selects between replica-centric and EC-centric layouts as TTL and pricing regimes shift.
\end{abstract}

\begin{IEEEkeywords}
SkyStore; geo-distributed storage; erasure coding; hybrid redundancy; mixed-integer linear programming; minimum-cost placement; cloud egress.
\end{IEEEkeywords}

%==============================================================================
\section{Introduction}
\label{sec:intro}

With the development of internet-scale services and electronic commerce, enterprises increasingly outsource data to geo-distributed cloud storage that spans multiple availability zones and pricing jurisdictions~\cite{skystore}.
Such deployments achieve flexibility for disaster recovery and locality-aware reads, yet \emph{redundancy management}---how many full replicas to keep versus when to stripe erasure-coded shards---directly couples storage invoices with cross-region bandwidth bills.

\textbf{Full-object replication} improves read latency and fault isolation at the expense of storage amplification and synchronization traffic.
\textbf{Erasure coding} trades lower stored bytes for coded fragments, placing stress on shard placement, reconstruction reads, and metadata bookkeeping.
Pure replica stacks and pure EC stacks rarely minimize total dollar cost simultaneously under realistic egress metering.

\textbf{Nevertheless}, widely used trace-driven frameworks center evaluation on replica-centric policies and TTL heuristics.
\textbf{SkyStore} exemplifies this line: it enables reproducible comparison of caching and replication choices against public or proprietary traces~\cite{skystore}.
\textbf{However}, integrating EC as a first-class redundancy mode---selected alongside replicas under one objective---remains under-articulated in that baseline narrative.

\textbf{Moreover}, when each region offers distinct storage and egress prices, enumerating ``second replica here'' versus ``two shards there'' combinations scales poorly; manual knob tuning cannot guarantee near-optimal layouts at trace time granularity.

\textbf{Therefore}, we propose \textbf{hybrid replica--EC redundancy guided by a MILP minimum-cost formulation}: binary decisions toggle EC versus extra replicas, while linear coefficients approximate storage duration (horizon for the pinned base replica, TTL-derived windows elsewhere) plus one-shot placement transfer proportional to object or shard volume.

The \textbf{main contributions} of this paper are summarized as follows:
\begin{itemize}
\item \textbf{Firstly}, we analyze limitations of replica-only decision narratives for geo-distributed cloud storage under heterogeneous pricing, and motivate coupling EC placement with replica placement inside one optimization shell (Section~\ref{sec:problem}).
\item \textbf{Secondly}, we formalize a compact MILP whose objective aggregates mandatory base storage, optional replica terms, and EC shard terms consistent with a \(2{+}1\) Reed--Solomon storyline without introducing a third parity decision variable; we spell out variables, linear constraints, and coefficient construction tied to trace snapshots (Section~\ref{sec:design}).
\item \textbf{Finally}, we prototype the policy inside a SkyStore-derived simulator, describe experimental methodology over trace workloads, and report how MILP-guided hybrid redundancy shifts optimal layouts relative to replica baselines (Sections~\ref{sec:eval}--\ref{sec:conclusion}).
\end{itemize}

The rest of the paper is organized as follows.
Section~\ref{sec:related} surveys related work on redundancy, erasure coding placement, and optimization-based storage control.
Section~\ref{sec:problem} formalizes the system model, workload assumptions, and the optimization gap relative to SkyStore-style evaluation.
Section~\ref{sec:design} presents the MILP-centric hybrid redundancy design and its interface to the simulator.
Section~\ref{sec:eval} describes experiments and metrics.
Section~\ref{sec:conclusion} concludes and outlines future work.

%==============================================================================
\section{Related Work}
\label{sec:related}

Related work sections in venue templates often open with scope framing~\cite{milp_survey_storage}; we follow that habit by partitioning references along redundancy mechanisms and optimization methodology.

\subsection{Geo-Distributed Replication and Trace-Driven Evaluation}
Multi-region replication improves durability and read locality; trace replay frameworks quantify dollar cost under realistic cloud price sheets.
SkyStore connects trace ingestion, TTL reasoning, and monetary accounting~\cite{skystore}.
\textbf{However}, these studies predominantly optimize replica counts or eviction policies rather than jointly selecting replica layouts \emph{and} erasure-coded shards under a unified MILP surrogate.

\subsection{Erasure Coding in Cloud and Wide-Area Storage}
Erasure coding from classical Reed--Solomon schemes to locally repairable codes trades storage for repair bandwidth~\cite{milp_survey_storage}.
Prior systems emphasize static stripe width or cluster-only placement.
\textbf{Nevertheless}, few works co-design EC shard placement with optional full replicas at object granularity for geo-distributed traces subject to per-edge egress billing.

\subsection{Optimization and MILP for Storage Resource Allocation}
Mixed-integer programs allocate shards, caches, or backup copies subject to capacity budgets~\cite{milp_survey_storage}.
Static solves over sliding windows amortize solver latency.
\textbf{Moreover}, integrating TTL snapshots into coefficients aligns MILP refreshes with telemetry buckets similar to histogram eviction policies.

\subsection{Positioning Against Prior Art}
Compared with replica-only SkyStore evaluations and with EC papers that ignore trace-driven regional pricing, we unify \textbf{hybrid redundancy selection} and \textbf{multi-region placement} under a \textbf{minimum-cost MILP}, grounded in the same simulator hooks used for accounting reproducibility.

%==============================================================================
\section{Problem Definition}
\label{sec:problem}

This section mirrors the analytical posture of Section~III in reference IEEE papers: articulate observations about incumbent schemes before detailing our formulation.

\subsection{System and Cost Model}
\textit{[Outline: define regions, directed pricing graph (storage node attributes, edge egress costs), logical objects, optional fixed-base pinning, trace-driven requests.]}
Let \(\mathcal{R}\) denote regions, each with storage price ladders mapped to simulation seconds.
Cross-region transfers incur charges proportional to bytes shipped along edges matching cloud egress conventions.

\subsection{Workload Assumptions}
\textit{[Outline: trace columns (operation, object key, issue region, timestamps, sizes); how TTL histograms approximate residency for non-base replicas/shards.]}
Highlight asymmetry between horizon-priced base replicas and TTL-priced ephemeral placements.

\subsection{Limitations of Baseline Heuristics}
Enumerate combinatorial growth when deciding EC versus an extra replica, and motivate MILP as a tractable surrogate when decision cardinality stays bounded (e.g., one extra replica, two outbound stripes).

\subsection{Optimization Goal}
State the minimum-cost objective at a decision epoch: mandatory base term \(+\sum_r s^{\mathrm{rep}}_r x_r + \sum_{r,k} s^{\mathrm{ec}}_{r,k} y_{r,k}\), referencing Section~\ref{sec:design} for coefficient semantics.

%==============================================================================
\section{Hybrid Redundancy Design and MILP Formulation}
\label{sec:design}

\subsection{Architecture Overview}
\textit{[Outline: SkyStore-derived simulator pipeline --- ingestion, RegionManager accounting, placement policy hook on write path, MILP refresh cadence tied to histogram buckets.]}

\subsection{Introducing Erasure Coding alongside Replicas}
Describe \(K_{\mathrm{data}}\), shard sizing \(S_{\mathrm{shard}} = S_{\mathrm{obj}}/K_{\mathrm{data}}\), parity folded into base replica narrative, and solver dimensions (\texttt{NUM\_EC\_SHARD\_SLOTS}=2).

\subsection{Decision Variables and Constraints}
Present binary \(z\) (EC mode), \(x_r\) (full replica at \(r\neq\)base), \(y_{r,k}\) (stripe \(k\) at \(r\)).
Summarize constraints: \(x_{\mathrm{base}}=0\) under FB; \(\sum_{r\neq \mathrm{base}} x_r \le 1\); \(\sum_r y_{r,k}=z\); \(\sum_k y_{r,k}\le 1\); coupling \(x_r+z\le 1\) on non-base rows.

\subsection{Coefficient Construction}
Explain horizon integration for base replica constants, TTL snapshots feeding \(s^{\mathrm{rep}}\) and \(s^{\mathrm{ec}}\), and \(\mathrm{edgeCost}(\mathrm{issue}, r)\times\)volume terms sourced from graph self-loops (cost zero) versus inter-region edges.

\subsection{Solver Integration}
Discuss Gurobi invocation, time limits, MIP gap, strict failure handling, and caching decisions per issue-region profile.

\subsection{Complexity and Practical Controls}
Windowing, warm-start considerations, and fallback policies when solver latency exceeds SLA.

%==============================================================================
\section{Experimental Evaluation}
\label{sec:eval}

\subsection{Implementation Environment}
\textit{[Outline: Java simulator modules (e.g., \texttt{PolicyMilpECnewV3}, \texttt{GurobiEc21MilpBuilder}), hardware, Gurobi version.]}

\subsection{Datasets and Traces}
Describe IBM Cloud Object Storage traces or others, region counts (2/6/9 configurations), preprocessing such as smoke subsets.

\subsection{Baselines}
Replica-only SkyStore policies, EC-disabled ablations, heuristic TTL matchers.

\subsection{Metrics}
Total simulated USD (storage + egress), MILP solve time, layout churn, sensitivity to TTL buckets.

\subsection{Results}
\textit{[Outline: reserve figures/tables --- cost savings vs.\ baselines, breakdown by component, discussion when MILP prefers EC versus second replica.]}

%==============================================================================
\section{Conclusion}
\label{sec:conclusion}

Cloud storage operators must reconcile replicas with erasure coding under heterogeneous pricing.
We extended a SkyStore-class evaluator with hybrid redundancy and MILP-based minimum-cost placement decisions.
Experiments confirm that joint optimization responds to TTL and egress shifts more systematically than replica-only heuristics.

\textbf{Future work} includes richer EC parameter families (\(k{+}m\) beyond \(2{+}1\)), online rounding between MILP epochs, and tighter coupling to repair traffic models.

%==============================================================================
\section*{Acknowledgment}
Placeholder acknowledgment text.

\bibliographystyle{IEEEtran}
\bibliography{refs}

\end{document}
